\begin{abstract}
Memes derive their communicative power not from semantic alignment between image and text, but from
\emph{intentional incongruity}—a smiling dog in a burning room, a deadpan face paired with absurdist
caption. Yet existing multimodal systems assume modalities agree and consequently collapse on the
conflict-dominant, non-literal language that defines meme affect.
We present \textbf{MemeSonic}, a unified framework for meme affective understanding and expressive
audio generation, built on a central premise: the \emph{incongruity} between image and text is not
noise to suppress but a differentiable signal to model.
To capture this signal, our framework first introduces independent multimodal encoding: image and text
are each mapped to their own emotion distributions, and a continuous cross-modal incongruity score is
derived from their divergence—making the tension between modalities an explicit, learnable feature
rather than an artifact to discard.
This incongruity score then drives \emph{Incongruity-Aware Mixture-of-Experts (MMOE) Fusion} for
affective classification. Across ten tasks on three meme benchmarks (MET-Meme, Memotion~7k,
MMSD~2.0), Incongruity Fusion achieves the best Macro-F1 on five tasks and MMOE achieves best on
four, outperforming standard fusion baselines by up to \textbf{+12.1\%} accuracy.
Beyond classification, we propagate the same incongruity signal into audio generation via an
\emph{Incongruity-to-Prosody Adapter}, enabling expressive speech synthesis that mirrors the
affective tension of the source meme. Human evaluation (n=84) shows MemeSonic achieves AMOS~3.95
vs.\ 1.68 for Vanilla TTS, with 84.5\% preference in A/B tests.
Finally, to assess whether audio carries affective signal beyond image and text, we conduct a
controlled tri-modal alignment experiment using Emotion2Vec embeddings, finding that naive
heterogeneous fusion yields zero modal contribution, while a learned linear projector bridges the
modality gap and raises sentiment accuracy from 27\% to 82\%—confirming that audio is a
non-redundant affective modality once properly aligned.
\end{abstract}
